\documentclass[11pt,a4paper]{article}
\usepackage[margin=2.2cm]{geometry}
\usepackage[T1]{fontenc}
\usepackage[utf8]{inputenc}
\usepackage{booktabs}
\usepackage{array}
\usepackage{hyperref}
\usepackage{enumitem}
\title{Sensor fusion hypothesis with MMC5983MA + IIS2ICLX \\ \small{v4-5reps simulated study}}
\author{MOOVIA experimental workflow}
\date{\today}
\begin{document}
\maketitle
\section{Context}
This document studies a hardware hypothesis over one real capture only: \\path{v4-5reps-moovia-session-1775503335652.json}.
The objective is not to claim validated hardware performance. The objective is to create a documented simulation package for the next hardware iteration.
\section{What each chip really returns}
\begin{itemize}[leftmargin=1.4em]
\item \textbf{MMC5983MA}: 3-axis magnetic field, optional temperature and data-ready/status. This is the sensor that can plausibly constrain yaw.
\item \textbf{IIS2ICLX}: 2-axis acceleration / inclinometer output, optional temperature, timestamp and FIFO/status. This helps roll, pitch and vertical reference, but not absolute yaw.
\item \textbf{Important limit}: the inclinometer does not provide heading on its own. Yaw observability still depends on the magnetometer or another absolute reference.
\end{itemize}
\section{Simulation setup}
\begin{itemize}[leftmargin=1.4em]
\item Base session: 5 real repetitions from the raw IMU export.
\item Magnetometer stream simulated at 200 Hz with world magnetic field, hard-iron bias, light noise and quantization.
\item Inclinometer stream simulated at 208 Hz with gravity-dominant X/Y tilt observation, small dynamic leakage, bias and quantization.
\item Raw files and derived files are stored separately on purpose.
\item The fusion pass is experimental and postprocessed. The production pipeline was not replaced.
\end{itemize}
\section{IMU-only vs hypothesis-level fusion}
\begin{tabular}{p{6cm}rr}
\toprule
Metric & IMU-only & Hypothesis fusion \\
\midrule
Rep count & 5 & 5 \\
Active-end height (m) & -11.271 & 0.003 \\
Active-end lateral (m) & 7.739 & 0.001 \\
Max lateral in window (m) & 7.739 & 0.092 \\
Mean local rep excursion (m) & n/a & 0.335 \\
Mean peak vertical speed (m/s) & 1.377 & 0.944 \\
\bottomrule
\end{tabular}
\section{Interpretation}
\begin{itemize}[leftmargin=1.4em]
\item Rep count stays consistent: yes.
\item Active-end height reduction: 100.0\%.
\item Active-end lateral reduction: 100.0\%.
\item Mean heading confidence used in the hypothesis: 1.00.
\item Mean tilt confidence used in the hypothesis: 1.00.
\end{itemize}
The simulation suggests that the extra observability could plausibly improve endpoint closure and lateral stability for the 5-rep session. The strongest expected gain comes from giving yaw an external reference and giving tilt a cleaner gravity reference.
\section{What improves and what still does not}
\subsection*{What plausibly improves}
\begin{itemize}[leftmargin=1.4em]
\item Yaw drift should decrease once the magnetometer constrains heading.
\item Roll/pitch and vertical closure should become more stable once the inclinometer constrains gravity.
\item Rep-local velocity interpretation can become more believable because the rep path closes better.
\end{itemize}
\subsection*{What this simulation does not prove}
\begin{itemize}[leftmargin=1.4em]
\item It does not prove real magnetic robustness in the presence of soft-iron or environmental disturbances.
\item It does not prove the final calibration workflow, mounting sensitivity or synchronization costs.
\item It does not validate production-ready absolute position. This remains a hypothesis-level, assisted postprocessing pass.
\end{itemize}
\section{Conclusion}
Yes: these extra sensors can be useful enough to justify the next hardware iteration. The plausible division of labor is clear: MMC5983MA for heading/yaw, IIS2ICLX for vertical and tilt. The folder generated for this study should therefore be treated as a pre-hardware decision package, not as a hardware validation result.
\end{document}